\subsection{Medication (5123034)}
Die Daten für die Medikamentenverwaltung basieren auf einem Datensatz von Hugging Face \cite{HuggingFace_Dataset} , der im JSONL-Format bereitgestellt wurde. Dieser enthielt jedoch überwiegend lediglich Informationen zu den Medikamentennamen sowie zur Darreichungsform (z.,B. Flüssigkeit oder Tablette). Angaben zu Dosierungen oder Wirkstoffen waren nicht enthalten.
Für die Abbildung von Krankheiten wurde ein weiterer Datensatz von Mendeley \cite{Mendeley_Dataset} verwendet. Dieser stellte Informationen zu Krankheitsnamen und zugehörigen Symptomen bereit. Allerdings fehlten auch hier Angaben zu geeigneten Behandlungsmethoden oder spezifischen Medikamenten.
\subsubsection{Methodik der Datenaufbereitung}
Zur Erstellung eines für das System geeigneten Datensatzes wurden beide Quellen mithilfe von Python-Skripten verarbeitet, bereinigt und miteinander in Beziehung gesetzt. Die Verarbeitung erfolgte in mehreren Schritten:
Zunächst wurden die Medikamentendaten bereinigt und vereinheitlicht. Dabei wurden inkonsistente Schreibweisen normalisiert sowie irrelevante oder fehlerhafte Einträge entfernt. Anschließend wurde versucht, aus den Medikamentennamen heuristisch mögliche Wirkstoffe zu extrahieren. Diese Extraktion erfolgte regelbasiert, da keine strukturierten Wirkstoffinformationen vorlagen.
Im nächsten Schritt wurden die Medikamente kategorisiert, beispielsweise in Tabletten, Flüssigkeiten oder andere Darreichungsformen. Diese Klassifikation bildete die Grundlage für spätere Konsistenzprüfungen innerhalb der Daten.
Darüber hinaus wurden jedem Medikament eindeutige Identifikatoren zugewiesen und initiale Bestandsmengen definiert, die den verfügbaren Vorrat im Krankenhaus modellieren.
Auf Basis dieser vorbereiteten Daten wurde eine separate Dosis-Tabelle generiert. Hierzu wurde eine Menge typischer Dosierungswerte und -einheiten erstellt, die in weiteren Schritten referenziert werden konnten.
Anschließend wurde die Medikationstabelle erzeugt. Die Zeiträume der Medikation (Start- und Enddatum) wurden pseudozufällig generiert, wobei sichergestellt wurde, dass das Enddatum stets nach dem Startdatum liegt und eine minimale Behandlungsdauer eingehalten wird. Zusätzlich wurde bei der Verknüpfung von Medikamenten und Dosierungen darauf geachtet, dass nur kompatible Kombinationen entstehen (z.,B. flüssige Medikamente mit ml-Angaben).
Abschließend wurde die Diagnosetabelle erstellt. Hierbei wurden Diagnosen Patienten zugeordnet und zusätzlich Referenzen zu behandelnden Ärzten integriert.
\subsubsection{Begründung der Vorgehensweise}
Die gewählte Vorgehensweise ergibt sich primär aus den Einschränkungen der verfügbaren Datensätze. Da keine vollständigen medizinischen Informationen (insbesondere Wirkstoffe, Dosierungen oder Therapieempfehlungen) vorlagen, war eine synthetische Erweiterung der Daten notwendig.
Die heuristische Extraktion von Wirkstoffen stellt dabei einen pragmatischen Ansatz dar, um zumindest eine teilweise inhaltliche Anreicherung der Daten zu erreichen. Auch wenn diese Methode keine vollständige Korrektheit garantiert, verbessert sie die Nutzbarkeit der Daten im Anwendungskontext erheblich.
Die Einführung von Kategorien für Medikamente dient der Sicherstellung semantischer Konsistenz innerhalb des Systems. Insbesondere die Kopplung kompatibler Dosierungseinheiten an passende Medikamententypen trägt zur Erhöhung der Datenqualität bei und verhindert unrealistische Kombinationen.
Die pseudozufällige Generierung von Medikationszeiträumen sowie Diagnosen wurde gewählt, um realitätsnahe Testdaten zu erzeugen. Die bewusste Entscheidung, mehr Diagnosen als Patienten zu generieren, ermöglicht die Simulation von Multimorbidität, die in realen medizinischen Szenarien häufig vorkommt.
Insgesamt erlaubt dieser Ansatz die Generierung eines konsistenten und hinreichend realitätsnahen Datensatzes, der für die Entwicklung und Evaluation des Systems geeignet ist, auch wenn die zugrunde liegenden Originaldaten nur eingeschränkt vollständig waren.